\documentclass[conference]{IEEEtran}
\usepackage{graphicx}
\usepackage{amsmath,amssymb}
\usepackage{booktabs}
\usepackage{url}
\usepackage{xcolor}
\usepackage[colorlinks=true,linkcolor=blue,citecolor=blue,urlcolor=blue]{hyperref}

% auto-generated by scripts/make_tables.py -- do not edit
\newcommand{\TunedTestAcc}{63.1\%}
\newcommand{\TunedTestMacroF}{0.342}
\newcommand{\TunedTestEvFone}{0.212}
\newcommand{\TunedLatencyMean}{0.39}
\newcommand{\BaseTestAcc}{32.0\%}
\newcommand{\KVEqCases}{3}
\newcommand{\SpecEqCases}{4}


\title{TrialMatch: Evidence-Grounded Patient--Trial Eligibility Screening\\
with Cache-Aware, Tiled and Speculative Transformer Inference}

\author{\IEEEauthorblockN{Abdullah Amin}
\IEEEauthorblockA{Generative AI, Semester 7\\
Individual assignment submission}}

\begin{document}
\maketitle

\begin{abstract}
Matching a patient narrative against the inclusion and exclusion criteria of a clinical
trial requires per-criterion decisions that are both correct and traceable to evidence in
the record. We present TrialMatch, a 370M-parameter decoder-only Transformer
(SmolLM2-360M-Instruct) fine-tuned with Low-Rank Adaptation on the TrialGPT Criterion
Annotations dataset under a mandated patient-grouped 37/8/8 split. The system predicts one
of six preserved expert categories per criterion---including \emph{not enough information}
and \emph{not applicable}---and cites the supporting note sentences by identifier. Beyond
task accuracy, we study inference itself: an explicitly owned key--value cache with
verified greedy equivalence against an uncached reference across input/output lengths; a
tiled attention kernel with online softmax validated numerically against a materialised
reference and PyTorch SDPA; and a from-scratch greedy speculative decoder over a
135M/360M draft--target pair whose emitted tokens provably match plain greedy decoding.
All components are exposed through a browser application with a live benchmark view.
Test-set results: \TunedTestAcc{} accuracy and \TunedTestMacroF{} macro-F1 after
fine-tuning (base model: \BaseTestAcc{}), with evidence-sentence F1 of
\TunedTestEvFone{} (micro). We report honest negative results where expected---notably
that a 135M draft cannot amortise its overhead against a 360M target at this scale---and
 analyse eight representative failure modes.
\end{abstract}

\begin{IEEEkeywords}
clinical decision support, patient--trial matching, LoRA, KV cache, online softmax,
speculative decoding, evidence grounding
\end{IEEEkeywords}

\section{Introduction and Problem Statement}
% TODO(final text): motivate per-criterion eligibility screening; state the two
% deliverables (screening quality + efficient-inference mechanisms); list contributions.

\section{Related Work}
% Cite TrialGPT (Jin et al.), LoRA, FlashAttention, online softmax, speculative
% decoding (Leviathan; Chen), SmolLM2, Qwen2.5, RoPE, GQA.
% Discuss: what TrialGPT does (GPT-4 prompting + retrieval), why a compact open-weight
% fine-tuned model with constrained decoding is a fair, reproducible alternative.

\section{Materials and Methods}
\subsection{Dataset and Split}
We use the prescribed \texttt{ncbi/TrialGPT-Criterion-Annotations} corpus. The GPT-4
prediction columns are discarded at load time and can never enter features or metrics.
Unique patient identifiers are sorted, shuffled with \texttt{numpy.random.default_rng(42)}
and assigned 37/8/8 to train/validation/test (602/206/206 annotations), with an executed
leakage check and a published manifest of patient IDs, counts and fingerprints.

\subsection{Model Selection}
% TODO: compare two candidates (SmolLM2-360M vs Qwen2.5-0.5B/1.5B): parameters,
% license (Apache-2.0), context window, tokenizer, memory; justify the pick.

\subsection{Input/Output Schema and Fine-Tuning}
Each example renders the note with its \emph{dataset-native} sentence numbering, the
criterion type and text; the target is a two-line answer \texttt{LABEL: ...} /
\texttt{EVIDENCE: [ids]}. Loss is computed on answer tokens only (prompt masked with
$-100$). At inference the label slot is restricted by a character trie over the labels
legal for the criterion type; unparseable output maps to \emph{not enough information}
rather than a positive decision. % hyperparameters table: r=16, alpha=32, lr=1e-3, ...

\subsection{Inference Mechanisms}
\subsubsection{KV cache} We own the cache as an explicit object (append, truncate, byte
accounting) and drive the checkpoint's layers manually: rotary positions continue at the
cached offset, prefill uses an offset-aware causal mask, decode steps are single-query
attends over the full cache. The uncached reference recomputes the whole prefix per token.
\subsubsection{Tiled attention} The tiled kernel walks key blocks with running maxima and
normalisers (online softmax), rescaling the accumulator by $\exp(m_{\text{old}}-m_{\text{new}})$.
\subsubsection{Speculative decoding} A 135M draft proposes $\gamma$ tokens from its own
cache; the target verifies all of them in one pass; the matching prefix is accepted, the
first mismatch is corrected with the target's argmax, and a bonus token follows a full
accept. Both caches are resynchronised at round boundaries; the label trie is applied to
the target's argmax during verification.

\subsection{Application Workflow}
% TODO: describe app screens: note+criteria input, per-criterion results with evidence,
% aggregated verdict rule, provenance-complete export, live benchmark tab, disclaimer.

\section{Experimental Setup and Metrics}
Common protocol: deterministic greedy decoding (temperature 0), identical prompts for base
vs.\ tuned comparisons, fixed held-out test untouched until final evaluation.
Classification: accuracy, macro-F1, per-class precision/recall/F1 and confusion matrices
overall and per criterion type. Evidence retrieval: micro and macro precision/recall/F1;
records with an empty gold set are excluded from the macro mean and reported separately
(an empty prediction on an empty gold set counts as agreement in the micro counts).
Benchmarks: 2 warm-ups + 5 repetitions, median, \texttt{torch.cuda.synchronize()} on GPU,
peak-memory capture per run.

\section{Results and Discussion}
\subsection{Eligibility and Evidence Quality}
\begin{table*}[t]
\centering
\caption{Criterion-level eligibility classification and evidence-sentence
retrieval on the shared 8/8 validation/test patient splits (206 annotations
each). Decoding: deterministic greedy with trie-constrained label slot.}
\label{tab:main}
\begin{tabular}{llccccc}
\toprule
Model & Split & Acc.\ (\%) & Macro-F1 & Ev.-F1 (micro) & Ev.-F1 (macro) & Invalid \\
\midrule
base & validation & 25.7 & 0.068 & 0.000 & 0.000 & 1.000 \\
base & test & 32.0 & 0.097 & 0.000 & 0.000 & 1.000 \\
tuned & validation & 65.0 & 0.278 & 0.184 & 0.113 & 0.000 \\
tuned & test & 63.1 & 0.342 & 0.212 & 0.198 & 0.000 \\
\bottomrule
\end{tabular}
\end{table*}

\begin{table}[t]
\centering
\caption{Per-class precision/recall/F1 of the fine-tuned model on the
held-out test patients (overall and by criterion type).}
\label{tab:perclass}
\resizebox{\columnwidth}{!}{%
\begin{tabular}{lccc}
\toprule
Scope & Precision & Recall & F1 \\
\midrule
overall -- included & 0.417 & 0.370 & 0.392 \\
overall -- not included & 0.000 & 0.000 & 0.000 \\
overall -- not excluded & 0.695 & 1.000 & 0.820 \\
overall -- not enough information & 0.569 & 0.439 & 0.496 \\
overall -- not applicable & 0.000 & 0.000 & 0.000 \\
inclusion -- included & 0.417 & 0.370 & 0.392 \\
inclusion -- not included & 0.000 & 0.000 & 0.000 \\
inclusion -- not enough information & 0.569 & 0.763 & 0.652 \\
inclusion -- not applicable & 0.000 & 0.000 & 0.000 \\
exclusion -- not excluded & 0.695 & 1.000 & 0.820 \\
exclusion -- not enough information & 0.000 & 0.000 & 0.000 \\
exclusion -- not applicable & 0.000 & 0.000 & 0.000 \\
\bottomrule
\end{tabular}}
\end{table}

% TODO: interpret the tables; insert figures/ cm_*.png + training curves; discuss
% inclusion vs exclusion asymmetry (label imbalance: 387/602 exclusion in train).

\subsection{Correctness of the Inference Mechanisms}
The cached decoder produces token-identical outputs to the uncached reference (and to HF
\texttt{generate}) across \KVEqCases{} input/output configurations; the speculative
decoder matches plain greedy on \SpecEqCases{} configurations for $\gamma\in\{2,4\}$.
A subtle offset bug (multi-token verification queries masked by local rather than
absolute positions) surfaced exactly through this suite and was fixed---evidence that
token-equivalence testing is not ceremonial. % TODO: table correctness_summary numbers

\subsection{Efficiency Measurements}
\begin{table}[t]
\centering
\caption{KV-cache vs.\ uncached full recompute (greedy, identical inputs;
device cuda, median of 5 after 2 warm-ups).}
\label{tab:cache}
\resizebox{\columnwidth}{!}{%
\begin{tabular}{rrcccc}
\toprule
$S$ & new & cached tok/s & uncached tok/s & speedup & equal \\
\midrule
128 & 16 & 28.63 & 19.47 & 1.5$\times$ & True \\
128 & 48 & 28.07 & 16.53 & 1.7$\times$ & True \\
512 & 16 & 28.58 & 5.87 & 4.9$\times$ & True \\
512 & 48 & 27.68 & 5.50 & 5.0$\times$ & True \\
1024 & 16 & 19.94 & 2.55 & 7.8$\times$ & True \\
1024 & 48 & 28.02 & 2.56 & 11.0$\times$ & True \\
\bottomrule
\end{tabular}}
\end{table}

\begin{table}[t]
\centering
\caption{Forward-pass latency of one attention layer (cuda, fp32).}
\label{tab:attention}
\begin{tabular}{rrrrr}
\toprule
$S_q$ & heads & reference (ms) & tiled (ms) & SDPA (ms) \\
\midrule
256 & 6 & 0.44 & 7.53 & 0.14 \\
512 & 6 & 0.56 & 22.84 & 0.29 \\
1024 & 8 & 1.92 & 64.48 & 0.73 \\
1024 & 8 & 3.58 & 64.73 & 1.25 \\
2048 & 6 & 5.73 & 249.97 & 1.89 \\
\bottomrule
\end{tabular}
\end{table}

\begin{table*}[t]
\centering
\caption{Greedy speculative decoding ($\gamma=4$) vs.\ plain greedy on the
same prompts; draft SmolLM2-135M, target SmolLM2-360M.}
\label{tab:spec}
\begin{tabular}{rrccccc}
\toprule
prompt & greedy tok/s & spec tok/s & speedup & acc.\ rate & rounds & equal \\
\midrule
180 & 19.41 & 10.14 & 0.52$\times$ & 0.64 & 7 & True \\
180 & 26.20 & 12.62 & 0.48$\times$ & 0.64 & 7 & True \\
180 & 26.52 & 12.85 & 0.48$\times$ & 0.64 & 7 & True \\
180 & 26.61 & 13.22 & 0.50$\times$ & 0.64 & 7 & True \\
180 & 26.58 & 9.81 & 0.37$\times$ & 0.64 & 7 & True \\
180 & 26.43 & 13.20 & 0.50$\times$ & 0.64 & 7 & True \\
\bottomrule
\end{tabular}
\end{table*}

% TODO: interpret: cache speedup vs S; why tiled lags SDPA (eager Python loop vs fused
% kernel, HBM round-trips); why the 135M draft cannot beat the 360M target (acceptance
% rate vs draft-step cost at this scale).

\subsection{Failure Analysis}
% TODO: walk through >=8 errors from results/errors_tuned_test.json: missing evidence,
% hallucinated evidence, insufficient-information failures, contradictory verdicts,
% parse failures, latency outliers. Never treat absent evidence as eligibility proof.

\section{Limitations}
% TODO: 360M scale; single-dataset; template explanations; verdict aggregation is
% rule-based and unlabeled; greedy-only decoding; educational scope.

\section{Conclusion}
% TODO: summary of what was demonstrated and the defence-ready takeaways.

\section*{AI-Use Appendix}
% TODO: declare assistant usage (code drafting, debugging, report editing), per course
% policy; all reused code tested and understood; the live defence covers details.

\bibliographystyle{IEEEtran}
\bibliography{references}
\end{document}
